% =====================================================================
%  FICHE 8 — THÈME 5 — NIVEAU MOYEN — CORRIGÉ
% =====================================================================
\documentclass[11pt,a4paper]{article}
\usepackage[utf8]{inputenc}
\usepackage[T1]{fontenc}
\usepackage{lmodern}
\usepackage[french]{babel}
\usepackage{amsmath,amssymb,mathtools}
\usepackage{icomma}
\usepackage{siunitx}
\sisetup{output-decimal-marker={,},per-mode=power,group-separator={\,},group-minimum-digits=5}
\usepackage[version=4]{mhchem}
\usepackage{xcolor}
\usepackage{tikz}
\usepackage[most]{tcolorbox}
\usepackage{enumitem}
\usepackage{array}
\usepackage{fancyhdr}
\usepackage[hidelinks]{hyperref}
\usetikzlibrary{arrows.meta,positioning,calc}
\usepackage[a4paper,margin=2.1cm,headheight=26pt,headsep=10pt]{geometry}
\usepackage{titlesec}

\definecolor{primary}{RGB}{150,98,162}
\definecolor{primarydark}{RGB}{92,52,110}
\definecolor{excol}{RGB}{0,135,90}
\definecolor{attncol}{RGB}{200,40,55}

\tcbset{boxbase/.style={enhanced,breakable,boxrule=0.9pt,arc=2.5pt,left=3mm,right=3mm,top=2.3mm,bottom=2.3mm,fonttitle=\bfseries,coltitle=white,toptitle=0.6mm,bottomtitle=0.6mm}}
\newtcolorbox[auto counter]{correction}[1][]{boxbase,colback=excol!4,colframe=excol,title={\textsc{Corrigé}~\thetcbcounter\ \ifx\relax#1\relax\relax\else\textmd{--- \itshape #1}\fi}}

\newcommand{\rep}[1]{\textcolor{primarydark}{\textbf{#1}}}
\newcommand{\oui}{\textcolor{excol}{\textbf{Oui}}}
\newcommand{\non}{\textcolor{attncol}{\textbf{Non}}}

\pagestyle{fancy}
\fancyhf{}
\fancyhead[L]{\small\textcolor{primarydark}{\textbf{Fiche 8} — Thème 5 : Équilibre chimique et Le Chatelier}}
\fancyhead[R]{\small\textcolor{primarydark}{\textsc{Moyen — corrigé}}}
\fancyfoot[C]{\small\thepage}
\renewcommand{\headrulewidth}{0.8pt}
\renewcommand{\headrule}{\hbox to\headwidth{\color{primary}\leaders\hrule height \headrulewidth\hfill}}
\setlength{\parindent}{0pt}\setlength{\parskip}{3pt}

\begin{document}

\begin{center}
\begin{tikzpicture}
\node[rectangle,rounded corners=7pt,fill=excol,text=white,minimum width=16cm,
      minimum height=1.1cm,align=center,inner sep=6pt]
  {{\large\bfseries Thème 5 — Équilibre chimique et Le Chatelier}\\[1pt]{\small Niveau \textsc{Moyen} — corrigé détaillé}};
\end{tikzpicture}
\end{center}

\begin{correction}[calculer $K$ à partir des concentrations]
\begin{enumerate}[label=\alph*),leftmargin=2em,topsep=2pt,itemsep=2pt]
  \item $K = \dfrac{[\ce{PCl3}]\,[\ce{Cl2}]}{[\ce{PCl5}]}$ (chaque coefficient vaut $1$).
  \item $K = \dfrac{1{,}40 \times 0{,}50}{0{,}40} = \dfrac{0{,}70}{0{,}40} = \rep{1{,}75}$.
  \item $K = 1{,}75 > 1$ : l'équilibre est \emph{légèrement} déplacé vers les \rep{produits}.
\end{enumerate}
\end{correction}

\begin{correction}[équilibre hétérogène en phase gazeuse]
\begin{enumerate}[label=\alph*),leftmargin=2em,topsep=2pt,itemsep=2pt]
  \item Les solides $\ce{FeO}$ et $\ce{Fe}$ n'apparaissent pas : $K = \dfrac{p_{\ce{CO2}}}{p_{\ce{CO}}} = 0{,}4$.
  \item Chaque molécule de $\ce{CO}$ consommée forme une molécule de $\ce{CO2}$ : le nombre total de moles
        de gaz est conservé, donc $p_{\ce{CO}} + p_{\ce{CO2}} = \SI{1}{atm}$.
  \item Avec $p_{\ce{CO2}}=x$ et $p_{\ce{CO}}=1-x$ :
        \[ \frac{x}{1-x} = 0{,}4 \;\Rightarrow\; x = 0{,}4(1-x) \;\Rightarrow\; 1{,}4\,x = 0{,}4 \;\Rightarrow\;
           x = \frac{0{,}4}{1{,}4} = \textcolor{primarydark}{\boldsymbol{\tfrac{2}{7}}\ \si{atm}} \approx \SI{0.29}{atm}. \]
\end{enumerate}
\end{correction}

\begin{correction}[réécrire l'équation change $K$]
\begin{enumerate}[label=\alph*),leftmargin=2em,topsep=2pt,itemsep=2pt]
  \item Réaction inverse : $K' = \dfrac{1}{K} = \dfrac{1}{2} = \rep{0{,}5}$.
  \item Coefficients $\times 3$ : $K' = K^3 = 2^3 = \rep{8}$.
  \item Coefficients $\times \tfrac12$ : $K' = K^{1/2} = \sqrt{2} \approx \rep{1{,}41}$.
\end{enumerate}
\end{correction}

\begin{correction}[procédé de contact — améliorer le rendement]
Réaction exothermique ; gauche $=3$ mol de gaz, droite $=2$ mol.
\begin{enumerate}[label=\alph*),leftmargin=2em,topsep=2pt,itemsep=2pt]
  \item Ajout d'$\ce{O2}$ (réactif) : consommé $\to$ droite. \oui
  \item Abaisser $T$ : la réaction est exothermique, baisser $T$ déplace vers les produits. \oui
  \item Augmenter $p$ : vers le côté à moins de gaz ($2<3$), soit la droite. \oui
  \item Catalyseur : n'agit que sur la cinétique, pas sur le rendement. \non
  \item Argon à volume constant : ne change pas les pressions partielles. \non
\end{enumerate}
\emph{Compromis industriel :} baisser $T$ améliore le rendement mais ralentit la réaction ; en pratique on
choisit une $T$ modérée \emph{plus} un catalyseur.
\end{correction}

\begin{correction}[décomposition endothermique en récipient fermé]
Endothermique ; gauche $=0$ mol de gaz, droite $=1$ mol ($\ce{CO2}$).
\begin{enumerate}[label=\alph*),leftmargin=2em,topsep=2pt,itemsep=2pt]
  \item (i) $T\nearrow$ : sens endothermique $\to$ droite. \oui \quad
        (ii) ouvrir : $[\ce{CO2}]\searrow$, le système en reforme $\to$ droite. \oui \quad
        (iii) ajouter $\ce{BaCO3}(s)$ : solide pur, aucun effet. \non \\
        (iv) $p\nearrow$ : vers le côté à moins de gaz (la gauche, $0$ mol) $\to$ gauche. \non \quad
        (v) ajouter $\ce{CO2}$ (produit) $\to$ gauche. \non
  \item Réaction endothermique : augmenter $T$ \rep{augmente} $K$.
\end{enumerate}
\end{correction}

\end{document}
